\subsection{Why the internal measurement is designed this way}
\label{sec:design-rationale}

Causal analysis of language models is normally conducted within a single model. A
task is fixed, a minimally corrupted variant of the input is constructed, and
internal states are transplanted between the two runs to locate where the corrupted
information is recoverable \citep{vig2020causal,meng2022locating}. Extending this
across models is not a matter of repeating the procedure. When two models produce
different outputs on the same input, any internal difference between them admits two
explanations that the measurement cannot separate: the models differ
architecturally, or the models are effectively performing different tasks. Comparing
their internals under that condition compares an unknown to an unknown.

The usual way around this is to select cases on which the models already agree. That
substitutes a selection bias for the confound: the comparison is then restricted to
whatever the models happen to find easy, which is not a controlled choice of task.

The untreated condition is equally unsuitable as a reference. Which models look
anomalous under an untreated task is a function of the task and of which models were
tested: the same model is an outlier in one fleet and unremarkable in another. Nor is
untreated variance attributable within a fixed fleet --- the difference between a
model that selects the intended action on three of ten samples and one that selects
it on none has no fixed standard against which either is measured, and cannot be
separated from sampling, routing or tokenisation. A reference condition has to be
established rather than observed.

We take the other route. Rather than selecting models that behave consistently, we
engineer consistency, and then measure. The specification establishes a common,
reproducible operational decision across otherwise divergent frozen models
(Section~\ref{sec:results}); that convergence is verified per model--condition before
any activation is inspected. It then serves as a controlled reference: the input, the
emitted action and the decision criterion are held fixed across the compared models
by construction, so remaining internal differences are attributable to the models
rather than to divergent task behaviour.

This design is available only because the behavioural result precedes the internal
one. It carries two conditions. The perturbations applied to the specification must
be identical across the compared models, which we ensure with token-aligned rewrites
of individual lines. And only models that reach the reference condition enter the
comparison; a model that does not converge under the specification is not a missing
data point but a separate outcome of the same procedure, and we report it as one.

\subsection{What is measured, and what is named}
\label{sec:design-naming}

Internal measurements are usually reported by naming an internal quantity --- a
direction, a head, a layer at which something is said to occur --- and supporting the
name with an intervention whose output effect is consistent with it. Output
consistency underdetermines the name: many internal facts produce the same change,
and the measured quantities vary continuously across layers, inputs, models and
numerical precision, so any point given a name is a point selected from a continuum.

We name nothing internal. The outcome variable is a verified behavioural transition:
a model that selects one action in ten of ten samples and, under the specification,
selects the other in ten of ten. The two actions are mutually exclusive and the task
declares which one satisfies the objective, so the meaning of the measured quantity
is supplied by the task rather than inferred from the model. This is the shape of the
decisions that language models are deployed to make --- approve or reject, escalate
or not, urgent or not --- in which the semantics of the output is fixed by the system
that consumes it.

The internal measurements then describe what carries that transition and at what
depth it becomes sufficient. They do not assert what any state represents. The margin
$M$ is a log-ratio between two token families and is reported as such; where its sign
and the emitted token disagree, we report the disagreement rather than resolving it
in favour of the margin.

One condition of this design should be stated. Declaring the two actions fixes their
meaning, but it fixes \emph{correctness} only for tasks that have an unambiguous
intended action. Where the correct action is contested, the same procedure establishes
reproducibility against a declared criterion without establishing that the criterion
is right.


\subsection{The unit of analysis}
\label{sec:design-unit}

An accuracy figure reported for a model on a task presumes that the task, rather than
the particular input used to pose it, is the object being measured. That presumption
requires behaviour to be invariant across inputs that express the same task. In the
conditions reported here it is not: one question, one model fleet, one sampling
regime, and a correct-response rate ranging from $7.6\%$ to $95.9\%$ according to what
accompanies the question. The text surrounding a question is not a neutral carrier
for it.

We therefore treat the exact input as part of the experimental condition rather than
as an implementation detail of the task. Each model--condition cell is evaluated on
one byte-identical input; that input is hashed and the hash recorded with every
execution; and outcomes are reported per model--condition rather than aggregated
across formulations. We make no claim about a model's accuracy on the task. We make
claims about a model's behaviour under a declared input.

The same constraint governs the internal measurement. A causal result obtained on one
input describes the computation over that input. Where the behaviour it is credited
with is itself input-dependent, attributing the result to the task rather than to the
input overstates it. We therefore anchor the internal measurement to a verified
behavioural transition between two declared inputs, and report the depth at which the
transplanted state becomes sufficient rather than a property of the task.

The scope of this observation should be stated. Our conditions differ by what is
added to the question, not by rewording alone, so what is demonstrated here is
sensitivity to surrounding structure rather than to paraphrase. Invariance under
task-preserving rewording is not tested and is not claimed.
